\section{Related work}
\paragraph{Memorisation and generalisation in in-context learning.} The transition from finite-pool to continuous-prior Bayes
as pretraining diversity grows is established for linear regression \citep{raventos2023,panwar2024}, with asymptotic theory
\citep{lu2025} and a rational-strategy account \citep{wurgaft2025}; for sequential data it has been studied with discrete
Markov chains and explicit memorising and generalising Bayes references \citep{gibson2026}. Our references play the same
roles for continuous-state dynamical systems, with an added context-only Bayes reference distinct from the population one.

\paragraph{Networks as Bayesian predictors.} Prior-fitted networks approximate the posterior predictive of their training
prior \citep{muller2022}. Recent theory defines a Bayes gap for in-context learners and shows that the uncertainty arising
from the task mixture vanishes quickly in the number of examples \citep{wakayama2026}, consistent with our finding that
population knowledge matters mainly at short contexts. In regression with exponentially distributed coefficients, a
transformer outperformed least-squares, ridge and $\ell_1$-based estimators for small to moderate context lengths, which
the authors read as implicit prior inference \citep{cao2026}. Point-prediction comparisons are fragile for telling
posteriors apart \citep{aswadi2026}; we score full predictive distributions with the log score \citep{gneiting2007}.

\paragraph{Dynamical systems.} In-context learning of noisy linear dynamical systems has approximation theory with
least-squares references \citep{cole2025}. Context-only strategies can rival foundation models on dynamical data
\citep{zhang2026parroting}. In a toy example in its appendix, an amortised forecaster trained on a Bayesian AR(1) prior beats
a per-series AR(1) fit ``since the Bayesian model knows a little more information'' \citep{koshelev2023}; that is the
population-advantage idea for one prior, without a finite-library or a separate context-only Bayes reference, a
separability analysis, or pre-registration.

\paragraph{Variance, design and pre-registration.} Seed variance and stratified inference over runs
\citep{agarwal2021,sellam2022} shape our hierarchical bootstrap; pre-registration in machine learning has precedents
\citep{bertinetto2021}. Asking where rival models' predictions differ before collecting data is close in spirit to classical
optimal design for model discrimination, which we have not reviewed here; we did not find, in the sources we read, a map of
the separability of Bayes references for in-context learners made before any network exists. (The working name of this
project collided with an unrelated method \citep{zhou2023dynaicl}, hence the title.)
